%!TEX root = ../template.tex
%%%%%%%%%%%%%%%%%%%%%%%%%%%%%%%%%%%%%%%%%%%%%%%%%%%%%%%%%%%%%%%%%%%%
%% abstract-pt.tex
%% NOVA thesis document file
%%
%% Abstract in Portuguese
%%%%%%%%%%%%%%%%%%%%%%%%%%%%%%%%%%%%%%%%%%%%%%%%%%%%%%%%%%%%%%%%%%%%

\ExplSyntaxOn
\cs_if_exist:NT \TitleOnHeader
  { \TitleOnHeader }
\ExplSyntaxOff

\typeout{NT FILE abstract-pt.tex}%

O paradigma da Indústria 4.0 assenta na convergência de \glspl{CPS} e \glspl{CPPS} para alcançar processos produtivos ágeis e resilientes. Neste ecossistema, os \glspl{DT} fornecem réplicas virtuais capazes de espelhar ativos físicos através de telemetria bidirecional em tempo real. No entanto, o ciclo de vida da engenharia de \glspl{DT} industriais enfrenta elevadas barreiras de implementação. O desenvolvimento convencional depende de software de simulação proprietário, modelos 3D \gls{CAD} fragmentados desprovidos de semântica cinemática e sensorial, e mapeamento manual exaustivo de endereços de \gls{IO}, tornando a reconfiguração frequente proibitivamente morosa e sujeita a erros.

Para superar estas limitações, esta dissertação propõe uma arquitetura aberta, modular e orientada a modelos para a instanciação procedimental e sincronização em tempo real de \glspl{DT} industriais. Desenvolvida sob a metodologia \gls{DSR}, a solução estrutura-se em quatro camadas desacopladas: digitalização e modelação 3D de baixo custo poligonal (Camada~0), aplicação desktop sem código para mapeamento sensorial via \textit{raycasting} 3D e empacotamento \gls{AASX} (Camada~1), registo dinâmico e repositório de pacotes baseado em microsserviços Eclipse BaSyx~v2 em Docker (Camada~2), e montagem procedimental de cenários, simulação física de contactos, execução em modo duplo (\gls{VC} e telemetria em direto) e compensação cinemática de desvio em Unity 3D (Camada~3).

A metodologia foi validada empiricamente numa bancada mecatrónica didática modular \gls{FMS} baseada em componentes \textit{fischertechnik}, controlada por um \gls{PLC} Industrial Shields com FreeRTOS, bem como através de testes de stress paramétricos ao longo de 800 execuções entre 5 e 1000 módulos transportadores. Os resultados experimentais comprovam a instanciação completa da bancada piloto em cerca de 2{,}8\,s, sustentando taxas médias de renderização superiores a 480\,\gls{FPS}. Os ensaios de escalabilidade confirmam complexidade quase linear até 250 módulos e execução estável em tempo real até 500 módulos (acima de 90\,\gls{FPS}), demonstrando a viabilidade da \gls{AAS} como norma unificadora para o gémeo digital industrial.

\keywords{
  Indústria 4.0 \and
  Sistemas Ciber-Físicos \and
  Simulação Procedimental \and
  Digital Twin \and
  Comissionamento Virtual \and
  Asset Administration Shell \and
}